\subsection{Assymptotic Gilbert-Varshamov Bound}
The following subsection is based on \cite{huffman}[Section 2.1, 2.8 and 2.10] and \cite{notes_on_alg_geom_codes}[Chapter 5]. We start by stating and proving a few technical results as well as the non asymptotic Gilbert bound, and introduce some new notation, along the way, before reaching the main result of this subsection, namely the Gilbert-Varshamov bound.

\begin{lemma}\label{lem:number_of_elements_in_sphere}
  Let $x \in \mathbb{F}_q^n$, then $\abs{\overline{B_{r}}(x)} = V_{q}(n, r)$ where $V_{q}(n, r) = \displaystyle\sum_{i=0}^r \binom{n}{i} (q - 1)^{i}$.
\end{lemma}
\begin{proof}
  Let $S_i(x) = \left\{x' \in \mathbb{F}_q^n \middle| \d(x, x') = i\right\}$ for $i = 0, 1, \ldots, r$. Then $S_0(x), S_1(x), \ldots, S_{r}(x)$ are all pairwise disjoint and $\overline{B_r}(x) = \bigcup_{i = 0}^r S_i(x)$. This gives
  \begin{equation}\label{eq:cardinality_of_disjoint_sets}
    \abs{\overline{B_r}(x)} = \sum_{i = 0}^r \abs{S_i(x)}
  \end{equation}
  Now $\abs{S_i(x)} = \binom{n}{i}(q - 1)^i$, since $x' \in S_{i}$ differ from $x$ in $i$ out of $n$ entries. The result follows immediately by combining this with Equation \eqref{eq:cardinality_of_disjoint_sets}.
\end{proof}

Moving forward we will use the notation $B_{q}(n, d)$ be the maximum number of codewords in a linear code, of length $n$ and minimum distance at least $d$.

\begin{proposition}\label{prop:covering_radius}
  If $\mathcal{C}$ is linear code over $\mathbb{F}_{q}$ of length $n$, with minimum distance $d$ and $B_q(n, d)$ codewords then:
  \begin{equation*}
    \mathbb{F}_q^n = \bigcup_{c \in \mathcal{C}} \overline{B_{d - 1}}(c)
  \end{equation*}
\end{proposition}
\begin{proof}\textcolor{red}{det her bevis skal du se bort fra.}
  Assume for the sake of contradiction that there exists at least one $y \in \mathbb{F}_q^{n}$, such that $y \not \in \bigcup_{c \in \mathcal{C}} \overline{B_{d-1}}(c)$. Now pick
  \begin{equation*}
    x = \underset{y \in \mathbb{F}_q^{n} \setminus \bigcup_{c \in \mathcal{C}} \overline{B_{d-1}(c)}}{\arg \min} \left( d(y, \mathcal{C})\right)
  \end{equation*}


  then $k x \not \in \bigcup_{c \in \mathcal{C}} \overline{B_{d-1}}(c)$ for all $k \in \mathbb{F}_{q}$ \textcolor{red}{\textbf{TODO}}. This in turn imply that $\mathcal{C}' = \mathcal{C} \cup \Span \left\{x\right\}$ is another linear code, with minimun distance $d$ and length $n$, however $\abs{\mathcal{C}'} > \abs{\mathcal{C}}$, which contradicts the fact that $\mathcal{C}$ had the maximum number of codewords.
\end{proof}

We are now able to state and prove the Gilbert Bound which we will use to prove the asymptotic Gilbert-Varshamov bound.
\begin{corollary}[Gilbert Bound]\label{cor:gilbert_bound}
  Suppose $V_{q}(n, r)$ is defined as in Lemma \ref{lem:number_of_elements_in_sphere} then:
  \begin{equation*}
    B_q(n, d) \geq \frac{q^{n}}{V_{q}(n, d - 1)}
  \end{equation*}
\end{corollary}

\begin{proof}
  Let $\mathcal{C}$ be a linear code of length $n$, with minimum distance $d$ and $B_q(n, d)$ codewords, then $\mathbb{F}_q^n = \bigcup_{c \in \mathcal{C}} \overline{B_{d - 1}}(c)$ by Proposition \ref{prop:covering_radius}, combining this with Lemma \ref{lem:number_of_elements_in_sphere}, we get
  \begin{equation*}
    q^{n} = \abs{\mathbb{F}_q^n} \leq \sum_{c \in \mathcal{C}} \abs{\overline{B_{q}}(c)} = B_q(n,d) V_q(n, d - 1). \qedhere
  \end{equation*}
\end{proof}

Continuing we define the following function which will used in the asymptotic Gilbert-Varshamov bound.
\begin{definition}
  The function $H_q: [0, (q - 1)/q] \to \mathbb{R}$ defined as $H_q(0) = 0$ and $H_q(x) = x \log_{q}(q-1) - x \log_{q}(x) - (1-x)\log_q(1 - x)$ when $x \neq 0$ is called the \textbf{$q$-ary entropy function}
\end{definition}
We will now state Lemma 2.10.3, from \cite{huffman}, but we will omit it's computation heavy proof.
\begin{lemma}\label{lem:entropy_limit}
  Let $0 \leq \delta \leq (q - 1) / q$ where $q \geq 2$, then $\lim_{n \to \infty} \frac{1}{n} \log_q(V_q(n, \floor{\delta n})) = H_q(\delta)$.
\end{lemma}


Moving on we introduce two invariants which can be used to gauge the quality of a code.
\begin{definition}
Let $\mathcal{C}$ be a $[n, k, d]_q$ code, then the invariants $R(\mathcal{C}):= k / n$ and $\delta(\mathcal{C}):= d / n$ is called the \textbf{transmission rate} and \textbf{relative distance} of $\mathcal{C}$ respectively.
\end{definition}
Following the definition a natural question arise, namely why would one consider the invariants $R(\mathcal{C})$ and $\delta(\mathcal{C})$ instead of using $n, k$ and $d$ directly, when gauging the quality of a $[n,k,d]_q$ code? To answer this consider repetition code of length $n$, as in example \ref{exmp:repetition_code}, when $n$ increases then so will $d$ since $d = n$, however $\delta(\mathcal{C})$ will stay constant as $n$ increases. Also as we are still only able to encode a single symbol we are sacrificing transmission rate.

As $B_q(n, d)$ is the maximum number of codewords in some linear code $\mathcal{C}$, with length $n$ and minimum distance at least $d$, which is a subspace of $\mathbb{F}_q^{n}$, we have $\dim(\mathcal{C}) = \log_q(B_q(n, d))$. Thus it makes sense to define what we shall reffer to as the \textit{maximal rate}:
\begin{equation*}
  R^{*}(n, d) = \frac{\log_{q}(B_{q}(n, d))}{n}.
\end{equation*}
We will now investigate the behavior of $R^{*}(n, d)$ as $n \to \infty$, for this we define, the function $\alpha_q^{lin}: [0, 1] \to [0, 1]$ as a function of a relative distance $\delta$, more specifically:
\begin{equation*}
  \alpha_q^{lin}(\delta) = \underset{n \to \infty}{\lim \sup} \;R^{*}(n, \delta n),
\end{equation*}
Finally we can state and prove the Assymptotic Gilbert-Varshamov\footnote{The Varshamov bound is another bound for the parameters of linear codes, which can be used to prove the theorem as well.} Bound.
\begin{theorem}[Assymptotic Gilbert-Varshamov Bound]
  Let $0 \leq \delta \leq (q - 1)/q$, then
  \begin{equation*}
  \alpha_q^{lin}(\delta) \geq 1 - H_q(\delta).
  \end{equation*}
\end{theorem}
\begin{proof}
  First note that $B_q(n, \delta n) = B_q(n, \ceil{\delta n})$, as $B_{q}(n, d)$ is the maximum number of codewords of a code of length $n$ with minimum distance at least $d$.
  % NOTE: Hvis du undre dig over hvor uligheden kommer fra så test med $\delta n \in \mathbb{Z}$.
  Thus we have
  \begin{align*}
\alpha_q^{lin}(\delta) = \underset{n \to \infty}{\lim \sup} \; R^{*}(\delta) &= \underset{n \to \infty}{\lim \sup} \; \frac{\log_q(B_q(n, \ceil{\delta n}))}{n} \\&\geq \underset{n \to \infty}{\lim \sup} \; \frac{\log_{q}\left(\frac{q^{n}}{V_{q}(n, \ceil{\delta n} - 1)}\right)}{n} \geq \underset{n \to \infty}{\lim \sup} \; 1 - \frac{\log_{q}({V_{q}(n, \floor{\delta n})}}{n}
  \end{align*}
  where the first inequality follows from the Gilbert bound (Corollary \ref{cor:gilbert_bound}) and the last inequality follows since $\ceil{\delta n} - 1 \leq \floor{\delta n}$ implies $\frac{q^{n}}{V_{q}(n, \ceil{\delta n} - 1)} \geq \frac{q^{n}}{V_{q}(n, \floor{\delta n})}$. The rest follows by applying Lemma \ref{lem:entropy_limit}, as
  \begin{equation*}
  \lim_{n \to \infty}\frac{\log_{q}({V_{q}(n, \floor{\delta n})}}{n} = H_q(\delta) \implies \underset{n \to \infty}{\lim \sup} \;\frac{\log_{q}({V_{q}(n, \floor{\delta n})}}{n} = H_q(\delta). \qedhere
  \end{equation*}
\end{proof}
